%% titre: Méthodes — Montrer qu'un ensemble est non borné
%% matiere: Optimisation
%% niveau: Licence 3
%% annee: 2026-2027
%% ordre: 6
%% description: Construire une famille de points admissibles dont la norme tend vers l'infini.

\section{Démarche}

\begin{methode}[Montrer qu'un ensemble $S\subset\R^n$ est non borné]
\begin{enumerate}[label=\textbf{Étape \arabic*.}]
    \item \textbf{Construire une famille de points} : choisir une famille paramétrée
    \[
    P_t=(x_{1,t},\ldots,x_{n,t})\in\R^n,
    \qquad t\geq 0.
    \]

    \item \textbf{Vérifier l'admissibilité} : montrer que
    \[
    P_t\in S,\qquad \forall t\geq 0.
    \]

    \item \textbf{Étudier la norme} : calculer
    \[
    \|P_t\|=\sqrt{x_{1,t}^2+\cdots+x_{n,t}^2}
    \]
    et vérifier que
    \[
    \lim_{t\to+\infty}\|P_t\|=+\infty.
    \]

    \item \textbf{Conclure} : puisque les points de $S$ peuvent s'éloigner arbitrairement de l'origine, on obtient
    \[
    \boxed{S\text{ est non borné}.}
    \]
\end{enumerate}
\end{methode}

\begin{attention}
\begin{itemize}
    \item Il suffit de construire \emph{une seule} famille de points admissibles dont la norme tend vers $+\infty$.
    \item Il faut impérativement vérifier que les points construits appartiennent à $S$ pour toutes les valeurs du paramètre considérées.
    \item Il n'est pas nécessaire de déterminer toute la géométrie de l'ensemble $S$.
    \item Une seule coordonnée qui tend vers l'infini en valeur absolue suffit pour établir la non-bornitude.
\end{itemize}
\end{attention}

\begin{exemple}
Considérons l'ensemble
\[
S=\left\{(x,y)\in\R^2:
\begin{array}{l}
x+y\leq 4,\\
x+3y\leq 9
\end{array}
\right\}.
\]

Choisissons la famille de points
\[
P_t=(-t,0),\qquad t\geq 0.
\]

Vérifions les contraintes :
\[
-t+0=-t\leq 4,
\]
\[
-t+3(0)=-t\leq 9.
\]

Ainsi, $P_t\in S$ pour tout $t\geq 0$.

D'autre part,
\[
\|P_t\|=\sqrt{(-t)^2+0^2}=t.
\]

Par conséquent,
\[
\lim_{t\to+\infty}\|P_t\|=+\infty.
\]

On conclut que
\[
\boxed{S\text{ est non borné}.}
\]
\end{exemple}
